% SPDX-License-Identifier: AGPL-3.0-or-later
% Preamble for the home page's live examples (content/_index.md).
% Character protrusion and font expansion, as in the PDF: the viewer
% protrudes and expands as the document says, and only then.
\usepackage{microtype}
\usepackage{amsthm}
\newtheorem{theorem}{Theorem}
\usepackage{tikz-cd}
\usepackage{reflowtex}
% Like a sectioning command, it leaves the paragraph after it unindented.
\makeatletter
\NewDocumentCommand\pagetitle{o m}{%
  \noindent\IfValueT{#1}{{\small\scshape #1}\par\smallskip\noindent}%
  {\LARGE\bfseries #2\par}\@afterindentfalse\@afterheading}
\makeatother
% A card's title: bold at the text size – the face the page's headings use,
% so no new style – with a little space before the text. \cardtitle* is for
% a card that opens with a display: the title stays the line the display
% follows, so TeX sets no empty line above the display and, the title being
% short, uses \abovedisplayshortskip. (After \par, a display would open a
% paragraph of its own, an empty line and the full \abovedisplayskip.) The
% space above that display is the \medskip a text card has after its title,
% whichever of the two skips TeX takes (an amsmath alignment always takes the
% full one, an equation the short one after a short line).
\NewDocumentCommand\cardtitle{s m}{\noindent{\bfseries #2}%
  \IfBooleanTF{#1}{\abovedisplayskip=\medskipamount \abovedisplayshortskip=\medskipamount}{\par\medskip}}
% Section headings ragged right: at a narrow measure a heading breaks onto
% several lines, and the standard classes would justify them. Only the
% heading's style changes, not its spacing or font.
\usepackage{etoolbox}
\makeatletter
\patchcmd{\@sect}{\begingroup #6}{\begingroup #6\raggedright}{}{}
\patchcmd{\@ssect}{\begingroup #4}{\begingroup #4\raggedright}{}{}
\makeatother
